% =====================================================================
%  Thème 4 — Équations trigonométriques — NIVEAU FACILE
%  Corrigé
% =====================================================================
\documentclass[11pt,a4paper]{article}
\usepackage[utf8]{inputenc}
\usepackage[T1]{fontenc}
\usepackage{lmodern}
\usepackage[french]{babel}
\usepackage{amsmath,amssymb,mathtools}
\usepackage{icomma}
\usepackage{xcolor}
\usepackage{colortbl}
\usepackage{tikz}
\usepackage[most]{tcolorbox}
\usepackage{enumitem}
\usepackage{array}
\usepackage[a4paper,margin=2.2cm]{geometry}
\usetikzlibrary{arrows.meta,calc}

\definecolor{primary}{RGB}{214,51,108}
\definecolor{primarydark}{RGB}{140,20,64}
\definecolor{excol}{RGB}{0,135,90}

\DeclareMathOperator{\tg}{tg}
\DeclareMathOperator{\cotg}{cotg}
\newcommand{\degr}{^{\circ}}
\newcommand{\Z}{\mathbb{Z}}

\newcounter{exo}
\newcommand{\exo}[1]{\medskip\refstepcounter{exo}%
  \noindent{\large\bfseries\color{primarydark}Exercice \theexo}%
  \ \textcolor{primary}{\textbullet}\ \textbf{#1}\par\nopagebreak\smallskip}
\newcommand{\rep}{\textcolor{excol}{$\blacktriangleright$}\ }

\setlength{\parindent}{0pt}
\pagestyle{empty}

\begin{document}

\begin{center}
{\LARGE\bfseries\color{primary}Thème 4 — Équations trigonométriques}\\[2pt]
{\color{primarydark}\textsc{Niveau Facile — Corrigé détaillé}}
\end{center}
\hrule height 1pt
\medskip

\exo{Une équation en cosinus}
\rep On reconnaît $\dfrac12=\cos\dfrac{\pi}{3}$, donc $\cos x=\cos\dfrac{\pi}{3}$.
L'équivalence du cosinus ($\cos a=\cos b\Leftrightarrow a=\pm b+2k\pi$) donne les deux familles
\[
x=\frac{\pi}{3}+2k\pi \qquad\text{ou}\qquad x=-\frac{\pi}{3}+2k\pi,\qquad k\in\Z.
\]
\[
\boxed{\,S=\left\{\tfrac{\pi}{3}+2k\pi\ ;\ -\tfrac{\pi}{3}+2k\pi \ \middle|\ k\in\Z\right\}.}
\]

\exo{Une équation en sinus}
\rep Ici $\dfrac{\sqrt2}{2}=\sin\dfrac{\pi}{4}$, donc $\sin x=\sin\dfrac{\pi}{4}$.
L'équivalence du sinus ($\sin a=\sin b\Leftrightarrow a=b+2k\pi$ ou $a=\pi-b+2k\pi$) donne
\[
x=\frac{\pi}{4}+2k\pi \qquad\text{ou}\qquad x=\pi-\frac{\pi}{4}+2k\pi=\frac{3\pi}{4}+2k\pi,\qquad k\in\Z.
\]
\[
\boxed{\,S=\left\{\tfrac{\pi}{4}+2k\pi\ ;\ \tfrac{3\pi}{4}+2k\pi \ \middle|\ k\in\Z\right\}.}
\]

\exo{Une équation en tangente}
\rep $1=\tg\dfrac{\pi}{4}$, donc $\tg x=\tg\dfrac{\pi}{4}$. La tangente est de \textbf{période $\pi$}
($\tg a=\tg b\Leftrightarrow a=b+k\pi$) : une seule famille, avec un $+k\pi$ (pas $+2k\pi$).
\[
x=\frac{\pi}{4}+k\pi,\qquad k\in\Z.
\qquad\Longrightarrow\qquad
\boxed{\,S=\left\{\tfrac{\pi}{4}+k\pi \ \middle|\ k\in\Z\right\}.}
\]

\exo{Un cosinus nul}
\rep $0=\cos\dfrac{\pi}{2}$, donc $\cos x=\pm\dfrac{\pi}{2}+2k\pi$. Les deux familles
$\dfrac{\pi}{2}+2k\pi$ et $-\dfrac{\pi}{2}+2k\pi$ se regroupent en une seule (les angles $\frac{\pi}{2}$
et $-\frac{\pi}{2}$ sont diamétralement opposés, espacés de $\pi$) :
\[
x=\frac{\pi}{2}+k\pi,\qquad k\in\Z.
\qquad\Longrightarrow\qquad
\boxed{\,S=\left\{\tfrac{\pi}{2}+k\pi \ \middle|\ k\in\Z\right\}.}
\]

\exo{Encore un sinus}
\rep $\dfrac{\sqrt3}{2}=\sin\dfrac{\pi}{3}$, donc $\sin x=\sin\dfrac{\pi}{3}$ :
\[
x=\frac{\pi}{3}+2k\pi \qquad\text{ou}\qquad x=\pi-\frac{\pi}{3}+2k\pi=\frac{2\pi}{3}+2k\pi,\qquad k\in\Z.
\]
\[
\boxed{\,S=\left\{\tfrac{\pi}{3}+2k\pi\ ;\ \tfrac{2\pi}{3}+2k\pi \ \middle|\ k\in\Z\right\}.}
\]

\end{document}
